\clearpage
\begingroup
\fancyhead[L]{\small POSTHUMOUS PAPERS AND CORRESPONDENCE}
\renewcommand{\headrulewidth}{0pt}
\par\medskip
\Needspace{8\baselineskip}
\pdfbookmark[2]{[3]}{pallas-IX-fragment3}
\begin{center}[3.]\end{center}
\begingroup
\small
\newcommand{\IXmars}{\text{\AstroSymbols ♂}}
\newcommand{\IXpallas}{\text{\AstroSymbols ⚴}}
\par\noindent\hspace*{\parindent}[Finally, multiplication by the factors in equations 2) gives the expressions for
\(\frac{\mathrm{d}\log\mathrm{hyp}\,n}{n\,\mathrm{d}t}\) and \(\frac{\mathrm{d}\varphi}{n\,\mathrm{d}t}\). The former, when integrated as on pages 523---524, gives
the perturbations of the logarithm of the semimajor axis and those of the first part of the epoch. The latter
gives those of \(\varphi\). The results were as follows.
\[
\begin{alignedat}{2}
\delta\log a&=1\sin(\IXmars-3\IXpallas+10^0\ 27'\ 49'')
&\mspace{4mu}\delta\!\int n\,\mathrm{d}t&=0\overset{\prime\prime}{,}18\sin(\IXmars-3\IXpallas+280^0\ 27'\ 49'')\\
&\quad+1\sin(\IXmars-2\IXpallas+202\ 20\ 18)
&&\enspace+0{,}16\sin(\IXmars-2\IXpallas+292\ 20\ 18)\\
&\quad+6\sin(\IXmars-\IXpallas+90\ 54\ 56)
&&\enspace+0{,}28\sin(\IXmars-\IXpallas+180\ 54\ 56)\\
&\quad+11\sin(2\IXmars-5\IXpallas+124\ 17\ 38)
&&\enspace+8{,}14\sin(2\IXmars-5\IXpallas+34\ 17\ 38)\\
&\quad+2\sin(2\IXmars-4\IXpallas+162\ 24\ 56)
&&\enspace+0{,}12\sin(2\IXmars-4\IXpallas+252\ 24\ 56)\\
&\quad+2\sin(3\IXmars-7\IXpallas+76\ 26\ 28)
&&\enspace+0{,}38\sin(3\IXmars-7\IXpallas+166\ 26\ 28)\\
&\quad+2\sin(4\IXmars-10\IXpallas+171\ 18\ 55)
&&\enspace+0{,}79\sin(4\IXmars-10\IXpallas+81\ 18\ 55)\\
&\quad+1\sin(5\IXmars-12\IXpallas+120\ 44\ 45)
&&\enspace+0{,}28\sin(5\IXmars-12\IXpallas+210\ 44\ 45).
\end{alignedat}
\]
\[
\begin{aligned}
\delta\varphi&=0\overset{\prime\prime}{,}15\sin(\IXmars+335^0\ 22'\ 13'')\\
&\quad+0{,}49\sin(2\IXmars-5\IXpallas+106\ 33\ 27).
\end{aligned}
\]
\begin{center}
\begin{tabular}{@{}lrl@{}}
Secular part of \(\delta\varphi\): & daily & \(-0\overset{\prime\prime}{,}000\,0059\)\\
& annually & \(-0\overset{\prime\prime}{,}00215\).\\
\end{tabular}
\end{center}
\par\noindent\hspace*{\parindent}\(\delta\log a\) is given in units of the 7th decimal place. All equations are included whose magnitudes
reach one such unit or\ \(0\overset{\prime\prime}{,}10\), respectively.]
\par\smallskip\noindent VII.\hfill 75\par
\par\endgroup

\par\medskip
\SourceLeaf{600}
\Needspace{8\baselineskip}
\pdfbookmark[2]{[4]}{pallas-IX-fragment4}
\begin{center}[4.]\end{center}
\begingroup
\small
\par\noindent\hspace*{\parindent}[The series into which the quantities \(\frac{x'-x}{\rho^3}\), \(\frac{y'-y}{\rho^3}\), \(\frac{z'-z}{\rho^3}\) were expanded above
converge so slowly that, despite the smallness of the individual perturbation equations, 24 positions of Mars
and as many of Pallas had to be used as the basis of the calculation. Another procedure is possible here.
We multiply the quantities to be expanded by a factor \(\lambda\)
chosen so that the expansions of the products \(\frac{x'-x}{\rho^3}\lambda\) etc. converge as well as possible,
while the expansion of \(\lambda^{-1}\), though slowly convergent, can readily be carried as far as desired using analytical
formulas. The products \(\frac{x'-x}{\rho^3}\lambda\) etc. can then easily be expanded by interpolation,
using only a few positions of Mars, and the expansion of \(\lambda^{-1}\) can be continued with
little effort as far as seems necessary. Multiplying the two expansions then gives
the desired expansion of \(\frac{x'-x}{\rho^3}\) etc. This method can also serve to check the expansions obtained
by interpolation alone.
\par\noindent\hspace*{\parindent}The factor \(\lambda\) was determined in the following way.]
\par\endgroup
\clearpage
\endgroup
